\title[IoT]{\large A \alert{simple}, lightweight, highly extensible, event driven, layer-based framework for network applications in the spirit of \alert{node}.js written entirely in \alert{C }++}

\lecture[SNode.C]{\href{https://github.com/SNodeC/snode.c}{\Huge\alert{SNode.C}}}{SNode.C}

\part{Classification}
\section{Motivation}

\begin{frame}{Motivation}{Classification}
   \begin{center}
     \alert{\Huge \color{fh-ooe-red} SNode.C}\\\scriptsize\copyright{} 2020, 2021, 2022, 2023 Volker Christian\\
     \href{https://github.com/SNodeC/snode.c}{\beamergotobutton{https://github.com/SNodeC/snode.c}} \href{https://www.gnu.org/licenses /lgpl-3.0.en.html}{\beamergotobutton{GNU Lesser General Public License v3.0}}
   \end{center}
   The development of \texttt{SNode.C} began during the first Corona lockdown in the summer semester of 2020 as part of the master's course \alert{Networked and
     Distributed Systems}.
   \begin{description}[\alert{Design decision}]
     \item [\alert{Goal}] To familiarize students with the basic \alert{techniques and concepts of network and web frameworks}.
     \item [\alert{role model}] Well-known \alert{express} framework from \alert{node.js}.
     \item [\alert{Design decision}] originally
     \begin{itemize}
       \item Single tasking, single threaded.
       \item Non blocking, event driven, multiplexed I/O.
       \item "`select"' as the only I/O multiplexer
       \item API as similar as possible (JS vs. C++) to the node.js-express API.
       \item \alert{server-side} only.
       \item Only \alert{IPv4}
     \end{itemize}
     \item [\alert{Current status}] Much more than initially planned
   \end{description}
\end{frame}
